%%%%%%%%%%  scan idx 297  |  folio 284  %%%%%%%%%%
\((1=\cl(\mathcal O_X))\), \emph{et l'}%
\emph{homomorphisme canonique} (2.6.3.2)
\[
 K^{\bullet}(X)_{\naif}\longrightarrow K^{\bullet}(X)
\]
\emph{est un homomorphisme d'anneaux unifères}.

Soit \(f:(X,\mathcal O_X)\to(Y,\mathcal O_Y)\) un morphisme de topos annelés en anneaux commutatifs. Alors la compatibilité du produit tensoriel avec les images inverses montre que \emph{l'homomorphisme}
\[
 f^*:K^{\bullet}(Y)\longrightarrow K^{\bullet}(X)
\]
(resp.
\[
 f^*:K^{\bullet}(Y)_{\naif}\longrightarrow K^{\bullet}(X)_{\naif}
\]
) \emph{est un homomorphisme d'anneaux unifères}.
\end{enumerate}

\textbf{2.8. Augmentation de \(K^{\bullet}\) vers \(H^0(\ ,\mathbb Z)\).} Soit \((X,\mathcal O_X)\) un topos localement annelé (I 2.15.1). Rappelons qu'on a défini (I 6.11 à 6.14) un homomorphisme d'anneaux (unifères)
\[
 i:H^0(X,\mathbb Z)\longrightarrow K^{\bullet}(X)_{\naif},
\]
munissant ainsi \(K^{\bullet}(X)_{\naif}\) (et par suite aussi \(K^{\bullet}(X)\) via l'homomorphisme canonique \(K^{\bullet}(X)_{\naif}\to K^{\bullet}(X)\)) d'une structure d'algèbre (à unité) sur \(H^0(X,\mathbb Z)\). De plus (loc.cit.) l'algèbre \(K^{\bullet}(X)_{\naif}\) (resp. \(K^{\bullet}(X)\)) est augmentée vers \(H^0(X,\mathbb Z)\) par le \emph{rang}, et l'homomorphisme canonique \(K^{\bullet}(X)_{\naif}\to K^{\bullet}(X)\) est un homomorphisme d'algèbres augmentées.

Soit \(f:(X,\mathcal O_X)\to(Y,\mathcal O_Y)\) un morphisme de topos localement annelés. Les homomorphismes \(f^*:K^{\bullet}(Y)\longrightarrow K^{\bullet}(X)\) et \(f^*:K^{\bullet}(Y)_{\naif}\longrightarrow K^{\bullet}(X)_{\naif}\) sont compatibles (I 6.12, 6.13) avec les structures d'algèbres augmentées précédentes, ce qui s'exprime par la commutativité du diagramme (où \(\Rg\)
désigne le rang) :
